% =====================================================================
% SCHOOL VERSION — Appendix: Detailed Intervention Implementations
% Path-free thesis version: methods, design choices, and findings only.
% =====================================================================

\chapter{Detailed Intervention Implementations}
\label{app:interventions}

This appendix documents the design and implementation logic of the
selection-repair interventions evaluated in Chapter~\ref{ch:experiments}.
The main chapters summarize their empirical outcomes; the purpose here is to
make the intervention logic sufficiently explicit for methodological review,
replication, and extension. Implementation code, frozen configurations, and
supporting artifacts are maintained in the project repository.\footnote{\url{https://github.com/OscarTsao/CultureDx}}
Internal file paths, development-branch names, and software class names are
intentionally omitted because they are implementation details rather than part
of the scientific argument.

\section{Evidence-Reasoning Pipeline}
\label{app:evidence}

The clinically explicit evidence-reasoning intervention decomposes the
transcript into four stages:
\emph{symptom extraction}, \emph{somatization mapping},
\emph{criterion matching}, and \emph{evidence-brief assembly}.

\subsection{Symptom Extraction}

An LLM extracts clinically relevant symptom spans from transcript turns.
Long inputs are bounded to approximately 14{,}000 characters while preserving
turn boundaries so that quoted evidence remains attributable to a speaker and
local conversational context. Separate Chinese- and English-language prompt
templates are used, but both request the same structured output: a concise set
of symptom spans, their local context, and a confidence estimate.

\subsection{Somatization Mapping}

Extracted symptom spans are mapped to somatic-presentation criterion groups.
This stage is intended to recover clinically relevant bodily complaints that
may be expressed indirectly or colloquially rather than in diagnostic
terminology. A fallback LLM prompt handles spans that cannot be mapped by the
primary rule-based pass. The output is a set of candidate somatic mappings,
not a diagnosis.

\subsection{Criterion Matching}

Each extracted span is compared with the ICD-10 criterion text for the active
candidate disorders. The principal lexical matcher uses token-overlap
similarity rather than a dense embedding score. Text is normalized using
Unicode NFKC normalization, and scope-aware negation handling reduces the
matching score when the evidence explicitly contradicts the criterion. In the
tested configuration, a fixed contradiction penalty of 0.20 is subtracted
from the fused score.

A separate dense-retrieval path represents criterion text and transcript spans
with sentence embeddings and caches the resulting vectors. Temporal features
are computed independently so that explicit duration or course information can
be represented without being conflated with lexical overlap.

\subsection{Evidence-Brief Assembly}

Matched evidence is assembled into one evidence brief per candidate disorder.
The brief includes the supporting or contradicting span, criterion identity,
matching score, confidence, and temporal information when available. Matches
below a minimum confidence of 0.1 are excluded. Briefs are cached so that
downstream reranking and diagnostic agents receive a stable evidence view.

\subsection{Finding and Interpretation}

The intervention did not yield a stable primary-selection gain. Explicit
decomposition introduced additional failure modes, including extraction
errors, incorrect negation handling, and noisy somatization mappings. Because
all evidence originates from the same single-session transcript, the pipeline
can reorganize observed evidence but cannot recover duration, prior-course,
substance, medical, or exclusion information that was never elicited. The
result therefore supports a measured conclusion: clinically explicit
decomposition can improve inspectability, but noisy extraction may offset its
potential information value.

\section{Confidence Calibration}
\label{app:calibration}

The confidence-calibration intervention applies post-hoc Platt scaling to the
diagnostician's verbalized confidence. Calibration is fitted under five-fold
cross-validation, avoiding the use of the same case both to fit and evaluate
its calibration transform.

Two complementary calibration measures are reported:

\begin{itemize}
  \item \textbf{Expected Calibration Error (ECE):} a ten-bin comparison
        between confidence and empirical correctness;
  \item \textbf{Area Under the Risk--Coverage Curve (AURC):} the error
        accumulated as increasingly low-confidence cases are retained.
\end{itemize}

Raw and calibrated scores are projected to the paper-parent taxonomy before
evaluation. Calibration did not improve the system's primary-selection
operating point and did not consistently improve ECE or AURC. This negative
result is important because the bottleneck is principally a ranking problem:
a monotonic calibration map may improve probability interpretation without
changing which candidate is ranked first.

A broader cross-dataset calibration utility was also developed for
leave-one-dataset-out analysis, but it is methodologically distinct from the
canonical intervention reported in the thesis and is not used to support the
main selection claim.

\section{Learned Re-ranking}
\label{app:ranker}

Two families of learned rerankers were evaluated: a structured feature
reranker and a dense similarity reranker.

\subsection{Structured Feature Reranker}

A gradient-boosted ranking model is trained with out-of-fold predictions using
four groups of features:

\begin{enumerate}
  \item \textbf{Checker features}: number of criteria met, number of required
        criteria, and average checker confidence;
  \item \textbf{Evidence features}: number of evidence spans and mean evidence
        confidence;
  \item \textbf{Somatization features}: number of somatic mappings and number
        of somatic criteria judged met;
  \item \textbf{Cross-disorder features}: confidence margin, original rank
        position, and number of logic-confirmed disorders.
\end{enumerate}

The model improves point accuracy on the development lineage but transfers
poorly across corpora. This indicates that the fitted decision surface captures
dataset-specific ranking regularities rather than a robust clinical selection
rule.

\subsection{Dense Similarity Reranker}

A dense similarity reranker compares the case representation with candidate
disorder definitions under dense, sparse, multi-vector, and hybrid retrieval
modes. Each similarity score is fused with the frozen HiED ranking signal under
a common anti-leakage protocol. None of the tested modes provides a stable
cross-corpus improvement, and the richer multi-vector mode produces lower
point estimates in both evaluated corpora.

\subsection{Lexical Transfer Control}

A TF--IDF plus logistic-regression selector recovers part of the in-domain
ranking headroom but collapses under corpus shift. Training the same model
family directly on the external corpus restores performance within that
corpus, showing that the failure is one of transfer rather than an intrinsic
inability of lexical models to fit diagnostic labels. The result nevertheless
demonstrates that the learned lexical shortcut is corpus-dependent and does
not provide a portable solution to primary selection.

\section{Selector Fine-tuning}
\label{app:finetune}

Selector fine-tuning uses the same candidate-level features as the learned
reranking probes but optimizes the selector directly for the gold primary. The
evaluation uses a four-cell train/test matrix:

\begin{enumerate}
  \item train on LingxiDiag and test on LingxiDiag;
  \item train on LingxiDiag and test on MDD-5k;
  \item train on MDD-5k and test on LingxiDiag;
  \item train on MDD-5k and test on MDD-5k.
\end{enumerate}

Fine-tuning improves in-corpus point accuracy but does not transfer reliably.
The result narrows the claim boundary: a fitted selector can exploit corpus-
specific supervision, but this does not demonstrate a general solution to the
selection bottleneck across language style, corpus construction, and class
distribution shifts.

\section{Self-Consistency and Confidence-Informed Aggregation}
\label{app:selfconsistency}

Test-time aggregation resamples the diagnostician at
$K\in\{3,5\}$ and combines the committed primary using two rules:

\begin{enumerate}
  \item \textbf{Self-consistency}: plain majority vote over sampled primaries;
  \item \textbf{Confidence-informed self-consistency}: a confidence-weighted
        vote using the diagnostician's verbalized confidence.
\end{enumerate}

The $K=3$ result is computed from the first three samples of the same
five-sample set used for $K=5$, preserving a nested comparison. Ties are
resolved deterministically. A single-sample temperature ladder at
temperatures 0.0, 0.3, and 0.7 provides a control for sampling diversity.

None of the tested aggregation rules significantly improves Top-1 over greedy
Direct-Answer. The highest point estimate remains effectively unchanged, while
higher-temperature sampling reduces accuracy. These results suggest that the
selection error is not simply caused by insufficient sampling diversity or a
hidden majority consensus among independently sampled diagnostic commitments.

Aggregation is applied to the diagnostician because it is the component that
commits the primary diagnosis. Checker-level aggregation remains a separate
future-work question.

\section{Pairwise Differential Reasoning}
\label{app:pairwise}

The pairwise differential intervention targets cases in which two or more
candidates are both well supported and have closely matched criterion scores.
A close cluster is formed when candidate support exceeds the pre-specified
threshold and adjacent met-ratio differences fall below the distance
threshold. A differential agent then compares only the candidates inside this
cluster and selects the locally preferred primary.

The intervention is evaluated in three distinct ways:

\begin{enumerate}
  \item an independently generated comparison run;
  \item a controlled replay using identical frozen upstream outputs;
  \item a forced-override variant in which the pairwise choice replaces the
        diagnostician's committed primary.
\end{enumerate}

On recoverable close-cluster subsets, the local differential selects the gold
candidate in 28.6\% of LingxiDiag cases and 24.0\% of MDD-5k cases. However,
the controlled replay changes the ranked top candidate without changing any
committed primary under the reference finalization policy. The higher point
estimate in the independent comparison run is therefore not causally
attributable to the pairwise stage. When the pairwise choice is forced to
override the diagnostician, overall Top-1 falls sharply. Pairwise reasoning is
thus informative as a local ambiguity probe but is not supported as the
deployed selector.

\section{Residual-Evidence Loop}
\label{app:residual}

The residual-evidence intervention iteratively feeds checker outputs and
unresolved evidence back to the diagnostician. Its intended cycle is:

\begin{enumerate}
  \item identify criteria that remain insufficient or contradictory;
  \item summarize residual evidence and unresolved distinctions;
  \item ask the diagnostician to reconsider the candidate set;
  \item repeat under a bounded number of rounds.
\end{enumerate}

This intervention was evaluated through one-off sandbox orchestration rather
than a fully canonicalized release implementation. The available analyses are
sufficient to support the reported negative result, but the intervention
should be reimplemented as a clean, versioned module before further extension.

Across the tested model-family probes, the loop exhibits unstable trade-offs:
some configurations recover secondary labels while over-emitting on
single-label controls, whereas others remain conservative but fail to recover
missing labels. The result supports the thesis-level conclusion that repeatedly
feeding residual evidence back into the same transcript-only reasoning path
does not reliably resolve the selection bottleneck.

\section{Methodological Reproducibility and Scope}
\label{app:successor}

The interventions in this appendix differ in maturity. Some are integrated
into the canonical evaluation pipeline, whereas others originated as bounded
development probes. For reproducibility, each quantitative claim should be
interpreted together with its frozen configuration, dataset fingerprint,
prediction artifact, and statistical definition.

The methodological lessons are nevertheless independent of repository layout:

\begin{itemize}
  \item evidence extraction must separate observed evidence from missing
        information;
  \item calibration must be evaluated separately from ranking;
  \item learned rerankers require out-of-corpus validation;
  \item resampling methods must preserve nested sample sets and deterministic
        aggregation rules;
  \item local reranking requires a controlled replay to establish causal
        attribution;
  \item sandbox interventions should be canonicalized before being treated as
        reusable system components.
\end{itemize}

The repository preserves the implementation history and frozen artifacts, but
the thesis treats those materials as reproducibility support rather than as
part of the scientific prose.
